\chapter{Cognitive Load and Algorithmic Triage}

\section{The Threat of Notification Fatigue}
In legacy systems, user attention is a commodity to be strip-mined. Applications utilize intermittent variable rewards (push notifications, badges, endless scrolls) to maximize screen time, leading to severe altruism fatigue and burnout (Scenario Sigma). In the Sovereign Stack, human attention is the scarcest thermodynamic resource. 

\section{Algorithmic Triage and the Centaur UI}
Layer 6 implements strict "Algorithmic Triage" to defend the citizen's cognitive load. The UI acts as a cryptographic shield. When the Layer 4 Orchestrator generates a swarm of intents (e.g., thousands of micro-bounties during the Ecological Shock of Scenario Upsilon), the interface does not push a thousand notifications to the citizen's device.

Instead, the UI employs localized LLMs to act as a "Centaur" agent. The agent digests the incoming data streams, evaluates them against the citizen's current Trust Ring capacities and biological fatigue limits, and suppresses 99\% of the noise. The UI only surfaces highly curated, actionable intents that require the citizen's explicit physical or cryptographic intervention.

\section{Friction as a Feature}
While Web2 design obsessed over "frictionless" interfaces to maximize impulsive transactions, Layer 6 introduces deliberate, engineered friction. For irreversible actions (like burning a Soulbound Token or authorizing a dangerous L2 actuator), the UI physically slows down. It requires sustained haptic confirmation, biometric scans, or explicit vocalized consent. The interface forces the human to pause, engaging System 2 cognitive processing, ensuring that automation never overrides human agency in high-stakes scenarios.

\section{Iterative Refinement: Temporal Interface Decay}
Scenario Lambda (Cultural Commons) introduced the necessity of intentional forgetting. Digital artifacts in Web2 are permanent, leading to digital hoarding and cognitive clutter. Layer 6 implements \textit{Temporal Interface Decay}. Visual UI elements representing cultural or ephemeral intents actively "age"—their contrast lowers, textures fray, and they eventually fade from the default AR and screen views as their cryptographic TTL (Time-To-Live) expires, mimicking the natural decay of physical memory and easing the citizen's cognitive load.

\section{Iterative Refinement: Mesh-Based Identity Verification UI}
During Scenario Zeta (Micro-Lodging), the system lacked a visual method to rapidly verify Trust Rings in person. The UI now includes \textit{Cryptographic Access Control Overlays}. When two citizens meet, their wearable devices exchange ZK-proofs via BLE, and the AR/Smartphone UI instantly renders a localized visual graph, highlighting their shortest mutual Trust Ring path. This provides immediate, intuitive social proof without exposing the entirety of either user's global ledger history.
